\section{直角三角形的边角关系} \label{sec:triangle}

在直角坐标相关的几何问题中, 直角三角形有着重要的应用.
在弹幕中我们关注``长度''和``方向''两类几何属性,
在直角三角形中对应``边''和``角''.
这一节我们将探索直角三角形的各个边与角之间有怎样的联系.

如图\ref{fig:triangle},
在直角三角形$OAP$中, 点$O,A,P$称为直角三角形的\emph{顶点},
构成直角的两条线段$OA,AP$称为\emph{直角边},
另一条线段$OP$称为\emph{斜边},
三个角 (这次不是任意角)
$\angle{O}$ (或$\angle{AOP}$),
$\angle{A}$ (或$\angle{OAP}$),
$\angle{P}$ (或$\angle{OPA}$)
称为三角形的\emph{内角}, 有两个锐角和一个直角.
由于三角形的内角和为180\degree, 总有
\[
  \angle{A} = 90\degree,\quad
  \angle{P} = 90\degree - \angle{O}.
\]
所以我们只考察一个内角$\angle{O}$即可.

设图\ref{fig:triangle}
中直角三角形的直角边长度为
$|OA|=x$, $|AP|=y$, 斜边长度为$|OP|=r$,
内角$\angle{O}=\theta$.
我们对直角三角形边角关系的探讨将围绕
$x,y,r,\theta$这四个量进行.

\begin{figure}[htbp]
  \centering
  \begin{tikzpicture}[scale=2]
    \coordinate (O) at (0,0);
    \coordinate (A) at (2,0);
    \coordinate (P) at (2,1);

    \draw[help lines,-latex] (-0.5,0) -- (2.5,0);
    \draw[help lines,-latex] (0,-0.5) -- (0,1.2);

    \draw (O) -- (A) -- (P) -- cycle;
    \node[dot,label={below left:$O$}] at (O) {};
    \node[dot,label={below:$A$}] at (A) {};
    \node[dot,label={right:$P$}] at (P) {};
    \node[above right] at (1,0) {$x$};
    \node[left] at (2,0.5) {$y$};
    \node[above] at (1,0.5) {$r$};

    \draw (4mm,0) arc (0:{atan(0.5)}:4mm)
    node[right,pos=0.7] {$\theta$};
  \end{tikzpicture}
  \caption{直角三角形}
  \label{fig:triangle}
\end{figure}

\subsection{勾股定理}

勾股定理描述了直角三角形三条边的长度之间的关系.
该定理指出, 已知直角三角形的任意两条边的长度,
可以计算出第三条的边的长度.

勾股定理的文字表述为:
直角三角形的\emph{两条直角边长度的平方和}
等于\emph{斜边长度的平方}.
具体到图\ref{fig:triangle}的直角三角形$OAP$,
勾股定理表述为 \[ x^2 + y^2 = r^2. \]

勾股定理最常见的应用场景是计算两点之间的距离.
在LuaSTG中, 我们可以用函数
$\Dist(x_1,y_1,x_2,y_2)$
计算点$(x_1,y_1),(x_2,y_2)$之间的距离,
而\Dist 函数的内部原理就是勾股定理.
如图\ref{fig:dist-func},
平面上有两点$A=(x_1,y_1)$, $B=(x_2,y_2)$,
作图示直角三角形$ABC$, $AC$沿水平方向, $BC$沿竖直方向.
由$A,B$坐标可得两直角边的长度
\[
  |AC| = |x_1 - x_2|,\
  |BC| = |y_1 - y_2|.
\]
然后由勾股定理即可求出
\begin{align*}
  |AB| & = \sqrt{|AC|^2 + |BC|^2}                \\
       & = \sqrt{(x_1 - x_2)^2 + (y_1 - y_2)^2}.
\end{align*}

\begin{figure}[htbp]
  \centering
  \begin{tikzpicture}[scale=2]
    \coordinate (A) at (0,0);
    \coordinate (B) at (-2,1);
    \coordinate (C) at (-2,0);
    \draw[fill=black!10] (A) -- (B) -- (C) -- cycle;

    \node[
      dot,label={below right:$A=(x_1,y_1)$}
    ] at (A) {};
    \node[
      dot,label={left:$B=(x_2,y_2)$}
    ] at (B) {};
    \node[
      dot,label={below left:$C$}
    ] at (C) {};

    \node[below] at (-1,0) {$|x_1 - x_2|$};
    \node[left] at (-2,0.5) {$|y_1 - y_2|$};
    \node[above right] at (-1,0.5) {$\Dist(x_1,y_1,x_2,y_2)$};
  \end{tikzpicture}
  \caption{两点距离}
  \label{fig:dist-func}
\end{figure}

为了方便表述, 本教程引入LuaSTG的\Dist 函数,
用$\Dist(x_1,x_2,y_1,y_2)$表示点
$(x_1,y_1)$, $(x_2,y_2)$之间的距离;
用$\Dist(x,y)$
\footnote{虽然LuaSTG的\Dist 函数其实不接受只输入两个数值}
表示点$(x,y)$与原点的距离,
即$\Dist(x,y) = \Dist(0,0,x,y)$.

\subsection{三角函数}

三角函数描述了方位角与坐标比例的关系,
角度为锐角时, 这也对应直角三角形的内角与三边长度比例的联系.

具体地, 对于图\ref{fig:triangle}的直角三角形$OAP$,
若固定$\theta$为一个定值, 我们会发现,
无论三条边长度$x,y,r$如何变化,
$x,y,r$始终某种固定的比例关系,
而这个比例由内角$\theta$决定.
设平面上有一点$P=(x_p,y_p)$,
保持射线$OP$的方向不变, 移动点$P$,
我们会发现$x_p,y_p$也有一种固定的比例关系,
由$OP$方向决定.

\begin{definition}[三角函数]
  如图\ref{fig:trig-func},
  单位圆 (圆心为原点, 半径为1的圆)
  上有一点$P=(x_p,y_p)$, 相对原点的方位角为
  $\angle{xOP}=\theta$.
  换言之, 点$P$的极坐标为$\braket{1,\theta}$.
  根据点$P$的坐标, 我们定义以下函数:

  \begin{enumerate}[label=(\alph*)]
    \item 正弦函数$\sin(\theta) = y_p$,
    \item 余弦函数$\cos(\theta) = x_p$,
    \item 正切函数$\tan(\theta) = y_p / x_p$.
  \end{enumerate}

  这三个函数合称三角函数.
  虽然它们的倒数也称为三角函数, 但我们用不到所以就不管了.
\end{definition}

\begin{figure}[htbp]
  \centering
  \begin{tikzpicture}[scale=2]
    \draw[-latex] (-1.3,0) -- (1.3,0)
    node[right] {$x$};
    \draw[-latex] (0,-1.3) -- (0,1.3)
    node[left] {$y$};
    \draw (0,0) circle (1)
    node[below right] {$O$};

    \draw (0,0) -- (240:1.3);
    \node[
      dot,label={[background]above:$P = (-0.5,-0.8660)$}
    ] at (240:1) {};
    \draw[->] (0.2,0) arc (0:240:0.2)
    node[above,pos=0.5,background] {$\theta=240\degree$};

    \node[right,draw,align=left] at (2,0) {
      $\cos\theta = -0.5$ \\
      $\sin\theta \approx -0.8660$ \\
      $\tan\theta \approx 1.7321$
    };
  \end{tikzpicture}
  \caption{三角函数}
  \label{fig:trig-func}
\end{figure}

我们在初中学习平面几何时也有一套三角函数定义,
是用直角三角形的边长比例定义的.
这两种定义是相通的.
在图\ref{fig:triangle}中 (我特意画了一个坐标系),
直角三角形$OAP$的三边长度$x,y,r$有如下的比例关系:
\[x : y : r = \cos\theta : \sin\theta : 1.\]

\subsection{反三角函数与Angle函数}

运用三角函数, 我们可以由角度信息推导长度和坐标信息.
那么我们如何反过来, 由长度和坐标信息推导角度信息呢?
这就要用到反三角函数.

不严谨地说, 我们可以把反三角函数看作三角函数的``逆向操作''.
假设我们问``60\degree 的余弦值 ($\cos$) 是多少'', 这是三角函数.
那么我们反过来问: 哪一个角的余弦值是0.5? 这就是反三角函数.
每个三角函数对应一个反三角函数,
对于三角函数$\sin$, $\cos$, $\tan$,
数学上把对应的反三角函数记作
$\arcsin$, $\arccos$, $\arctan$,
而在编程中这些反三角函数命名为
\raw{asin}, \raw{acos}, \raw{atan}.

但是这个``逆向操作''需要一些特殊的处理才能成立.
当我们问``哪一个角的$\cos$值是0.5?'',
我们会发现答案有无数个:
$60\degree, -60\degree$, 以及它们的所有等价角.
作为一个函数必须要有\emph{确定性},
即, 接受相同的输入时, 必须有相同的输出
\footnote{
  编程上的函数没有这个限制, 比如随机数生成函数,
  但数学上一般是这样要求的, 这样才能进行确定的计算.
}.
为了让``反三角函数''具有确定性,
数学家人为地限制了每个反三角函数的输出范围.

以$\arccos$为例, 如图\ref{subfig:arccos}
要找一个角$\theta$使它的余弦值$\cos\theta$为$x_0$,
我们只需要在单位圆上找到一个点使得它的横坐标为$x_0$,
这样的点可能有两个, 我们取$x$轴上方的点,
把它相对于原点的方位角作为$\arccos x_0$的输出.

于是我们规定, $\arccos$的值域\footnote{输出数值的范围}为
$0\degree \le \arccos x_0 \le 180\degree$,
此时对于每一个$-1 \le x_0 \le 1$,
$\arccos x_0$都会输出唯一的角度.

类似地, 如图\ref{subfig:arcsin},
要计算$\arcsin{y_0}$ ($-1 \le y_0 \le 1$),
我们在单位圆上找一个点使得它的纵坐标为$y_0$,
这样的点也可能有两个, 我们取$y$轴右侧的点相对原点的方位角
作为$\arcsin{y_0}$的输出.
我们规定, $\arcsin$的值域为
$-90\degree \le \arcsin{y_0} \le 90\degree$.

如图\ref{subfig:arctan},
在平面上取一点$(1,k)$, 该点相对原点的方位角$\theta$
满足$\tan\theta = k$.
规定取$-90\degree < \theta < 90\degree$
的角度作为$\arctan{k}$的输出.

\begin{figure}[htbp]
  \centering
  \begin{subfigure}{0.4\textwidth}
    \centering
    \begin{tikzpicture}[scale=2]
      \clip (-1.7,-0.3) rectangle (1.7,1.7);
      \foreach \x in {-0.9,-0.6,...,0.9} {
          \draw[help lines] (\x,0) -- (\x,1.2);
          \draw[help lines] (0,0) -- (\x,{sqrt(1-\x*\x)});
        }

      \def\x{-0.6};
      \fill[gray!50] (0,0) -- (0.15,0)
      arc (0:{acos(\x)}:0.15) -- cycle;

      \draw[-latex] (-1.5,0) -- (1.5,0) node[right] {$x$};
      \draw[-latex] (0,-1.5) -- (0,1.5) node[left] {$y$};
      \draw (0,0) circle (1) node[below left] {$O$};

      \draw (\x,0) node[below] {$x_0$}
      -- (\x,1.2) node[above] {$x = x_0$};
      \draw (0,0) -- (\x,{sqrt(1-\x*\x)}) node[dot] {};
      \draw (0.15,0) arc (0:{acos(\x)}:0.15)
      node[above right,background,pos=0.5,outer sep=1pt]
        {$\theta=\arccos{x_0}$};
    \end{tikzpicture}
    \caption{$\arccos$函数}
    \label{subfig:arccos}
  \end{subfigure}
  \begin{subfigure}{0.28\textwidth}
    \centering
    \begin{tikzpicture}[scale=2]
      \clip (-0.3,-1.7) rectangle (1.7,1.7);
      \foreach \x in {-0.9,-0.6,...,0.9} {
          \draw[help lines] (0,\x) -- (1.2,\x);
          \draw[help lines] (0,0) -- ({sqrt(1-\x*\x)},\x);
        }

      \def\x{-0.6};
      \fill[gray!50] (0,0) -- (0.15,0)
      arc (0:{asin(\x)}:0.15) -- cycle;

      \draw[-latex] (-1.5,0) -- (1.5,0) node[right] {$x$};
      \draw[-latex] (0,-1.5) -- (0,1.5) node[left] {$y$};
      \draw (0,0) circle (1) node[below left] {$O$};

      \draw (0,\x) node[left] {$y_0$}
      -- (1.2,\x) node[below] {$y = y_0$};
      \draw (0,0) -- ({sqrt(1-\x*\x)},\x) node[dot] {};
      \draw (0.15,0) arc (0:{asin(\x)}:0.15)
      node[right,background] at (0.3,-0.15)
      {$\theta=\arcsin{y_0}$};
    \end{tikzpicture}
    \caption{$\arcsin$函数}
    \label{subfig:arcsin}
  \end{subfigure}
  \begin{subfigure}{0.28\textwidth}
    \centering
    \begin{tikzpicture}[scale=2]
      \clip (-0.3,-1.7) rectangle (1.7,1.7);
      \foreach \x in {-1.2,-0.9,...,1.2001} {
          \draw[help lines] (0,0) -- (1,\x);
        }
      \draw[help lines] (1,-1.4) -- (1,1.4);

      \def\x{0.9};
      \fill[gray!50] (0,0) -- (0.15,0)
      arc (0:{atan(\x)}:0.15) -- cycle;

      \draw[-latex] (-1.5,0) -- (1.5,0) node[right] {$x$};
      \draw[-latex] (0,-1.5) -- (0,1.5) node[left] {$y$};
      \draw (0,0) circle (1) node[below left] {$O$};

      \draw (0,0) -- (1,\x)
      node[dot,label={right:$(1,k)$}] {};
      \draw (0.15,0) arc (0:{atan(\x)}:0.15)
      node[right,background] at (0.25,0.15)
      {$\theta = \arctan{k}$};
    \end{tikzpicture}
    \caption{$\arctan$函数}
    \label{subfig:arctan}
  \end{subfigure}
  \caption{反三角函数}
\end{figure}

\subsubsection{反三角函数的值域问题}

在弹幕制作中使用反三角函数主要是为了根据坐标求解角度,
由于任意角的特性, 只要求出的角度范围能覆盖连续的360\degree 范围,
就可以满足一般的应用需求.

但是, 反三角函数只能覆盖180\degree 的范围,
这就导致我们直接用反三角函数求解角度时容易出错.
例如, 假设我们用$\arctan(y/x)$求点$(x,y)$相对原点的方位角,
那么当点$(x,y)$在$y$轴左侧时就会得到错误的结果
(甚至当该点在$y$轴上时无法计算$y/x$).

为了能够正确计算角度, 我们需要用$x,y$两个数值预先确定角度范围,
用这个角度范围来修正反三角函数的输出结果.
比如各编程语言提供的\raw{atan2}函数,
它的大致原理是这样的:
\[
  \raw{atan2}(y,x)=
  \begin{cases}
    \arctan(y/x),            & x>0,       \\
    \arctan(y/x)+180\degree, & y\ge0,x<0, \\
    \arctan(y/x)-180\degree, & y<0,x<0,   \\
    90\degree,               & y>0,x=0,   \\
    -90\degree,              & y<0,x=0,   \\
    未定义,                  & y=0,x=0.
  \end{cases}
\]

为了方便表述, 我们引入LuaSTG的\Angle 函数
\footnote{
  \Angle 函数的值域为
  $-180\degree < \theta \le 180\degree$.
},
用$\Angle(x_1,y_1,x_2,y_2)$
表示点$(x_2,y_2)$相对$(x_1,y_1)$的方位角;
另外定义, $\Angle(x,y)$
\footnote{虽然LuaSTG的\Angle 函数不能只输入两个数值.}
表示点$(x,y)$相对原点的方位角,
即$\Angle(x,y) = \Angle(0,0,x,y)$.

\subsection{直角坐标与极坐标的转换}

如图\ref{fig:default-coord},
在平面上同时建立直角坐标系和极坐标系,
直角坐标系和极坐标系的原点和极点重合,
单位长度相同, $x$轴正方向为极轴方向.
平面上的每一个点同时具有直角坐标和极坐标.

\begin{figure}[htbp]
  \centering
  \begin{tikzpicture}[scale=0.8]
    \def\len{3.7}
    \draw[-latex,very thick] (-\len,0) -- (\len,0)
    node[right] {$x$};
    \draw[-latex,very thick] (0,-\len) -- (0,\len)
    node[left] {$y$};

    \begin{scope}[help lines]
      \draw (45:3.5) node[right] {$45\degree$}
      -- (-135:3.5) node[left] {$-135\degree$};
      \draw (-45:3.5) node[right] {$-45\degree$}
      -- (135:3.5) node[left] {$135\degree$};
      \foreach \r in {1,2,3}
      \draw (0,0) circle (\r);
    \end{scope}

    \foreach \x in {-3,-2,-1,1,2,3} {
        \node[below left] at (\x,0) {$\x$};
        \node[below left] at (0,\x) {$\x$};
      }
    \node[below left,background] at (0,0) {$O$};

    \node[dot] at (1.5,1) {};
    \node[background,right] at (1.5,1)
    {$P = (x,y) = \braket{r,\theta}$};
  \end{tikzpicture}
  \caption{默认直角坐标系与极坐标系}
  \label{fig:default-coord}
\end{figure}

\subsubsection{极坐标 \textrightarrow 直角坐标}

设平面上一点$P$的直角坐标为$(x,y)$,
极坐标为$\braket{r,\theta}$.
根据三角函数的定义, 如果$r = 1$,
那么有$x = \cos\theta, y = \sin\theta$.
根据点$P$到原点的距离$r$做等比放缩,
我们容易得到将极坐标转换为直角坐标的公式:
\begin{align*}
  x & = r \cos\theta, \\
  y & = r \sin\theta.
\end{align*}

\subsubsection{直角坐标 \textrightarrow 极坐标}

反过来的转换则要复杂一些.
用勾股定理可以求解极径$r = \sqrt{x^2 + y^2}$.
在弹幕制作中, 直接调用$\Dist(0,0,x,y)$是更实际的选择.
至于极角$\theta$, 我们目前知道如果不为原点,
根据直角坐标$(x,y)$可以确定极角的一个等价取值.
有些弹设会考虑用反三角函数进行求解,
但是由于反三角函数的值域问题,
直接调用LuaSTG的\Angle 函数是一个更方便更正确的做法.
也就是:
\begin{align*}
  r      & = \Dist(x,y) = \sqrt{x^2 + y^2}, \\
  \theta & = \Angle(x,y).
\end{align*}

\subsection*{习题}

\begin{figure}[htbp]
  \centering
  \begin{tikzpicture}[scale=1]
    \draw (0,0) node[
      dot,label={left:$O$}
    ]{} circle (1);
    \foreach \a in {20,40,...,360}
    \node[dot] at (\a:1) {};
    \draw[->] (0,0) -- (40:1)
    node[above right]{$P$};
  \end{tikzpicture}
  \caption{龙纹弹示意图}
  \label{fig:long-wen-dan}
\end{figure}

\begin{enumerate}
  \item \comment{龙纹弹}
        龙纹弹也算是新人弹设的热门复刻项目了.
        为了让子弹始终围成一个圆形, zun有自己的谜之方法,
        而我们只需要每帧进行一个坐标的同步.
        对我们而言, 圆心的运动是简单的,
        设置一个隐形obj作为圆心进行直线运动,
        圆的半径等参数也可以由这个隐形obj更新.
        圆上的每个子弹读取圆心obj提供的圆心坐标, 半径等参数,
        更新自己的坐标, 从而保持自己始终在圆上.

        如图\ref{fig:long-wen-dan},
        设有一个子弹$P$, 在某一帧,
        它对应的圆心坐标为$(x_0,y_0)$,
        圆的半径为$R$, 该子弹相对于圆心的方位角为$\alpha$,
        求该帧子弹$P$的坐标.

  \item 证明: $(\cos\theta)^2 + (\sin\theta)^2 = 1$.
        (提示: 三角函数的定义)

  \item \comment{旋转公式v1*}
        设平面上有一点$(x,y)$, 将该点绕原点旋转$\phi$角
        (以逆时针为正方向), 得到的点的坐标$(x',y')$怎么表示?
        这就是二维旋转公式要解决的问题,
        它的标准形式为:
        \begin{align*}
          x' & = x \cos\phi - y \sin\phi, \\
          y' & = x \sin\phi + y \cos\phi.
        \end{align*}
        我个人比较讨厌这个公式, 因为这个公式很容易记错,
        很多弹设做的弹幕效果不对在群里提问,
        我们看了半天结果是旋转公式写错了, 非常折磨.
        最好是能够自己写一个lua函数实现这个旋转公式,
        需要旋转的时候统一调用函数, 以防止偶然写错.

        在\ref{sec:vector-base}节的习题会探讨二维旋转公式的
        另一种形式, 和我一样不想记这个标准形式的可以看一看.

        这一次的推导基于极坐标直角坐标转换和三角函数的和差角公式.
        设旋转前点$(x,y)$的极坐标为$\braket{r,\theta}$,
        显然, 旋转后的极坐标为$\braket{r,\theta+\phi}$.
        请你根据以下的两条公式, 补全二维旋转公式的推导过程:
        \begin{align*}
          \cos(\theta+\phi) & =
          \cos\theta\cos\phi - \sin\theta\sin\phi, \\
          \sin(\theta+\phi) & =
          \cos\theta\sin\phi + \sin\theta\cos\phi.
        \end{align*}

  \item \comment{方形造形术**}
        方形造形术是一个经典的异形开花弹案例.
        一个基本的异形开花弹, 各个子弹有相同的发弹点, 且都做匀速直线运动.
        这样的异形开花弹保持形状的背后思想是,
        把子弹的速度当作平面上的点来看待.
        如图\ref{fig:square-danmaku},
        当各子弹的速度对应点构成方形时, 异形开花弹就会保持方形.

        \begin{figure}[htbp]
          \centering
          \begin{tikzpicture}[scale=1.5]
            \draw[-latex] (-1.5,0) -- (1.5,0)
            node[right] {$x$};
            \draw[-latex] (0,-1.5) -- (0,1.5)
            node[left] {$y$};
            \node[below left] at (0,0) {$O$};

            \def\a{15};

            \begin{scope}[scale=1,rotate=\a]
              \draw[help lines] (-1,-1) rectangle (1,1);
              \node[left] at (-1,1) {$C$};
              \node[left] at (-1,-1) {$D$};
              \foreach \i in {0,1,2,5,8} {
                  \node[dot] at (1,-1+\i*2/8) {};
                }
              \foreach \i in {0,1,8} {
                  \node[right] at (1,-1+\i*2/8) {$B_\i$};
                }
              \draw[->] (0,0) -- (1.5,0)
              node[right] {$\alpha$};
              \draw[->,scale=0.95] (0,0) -- (1,-1+5*2/8)
              node[right] {$B_i$}
              node[above,pos=0.5] {$v_i$};
            \end{scope}
          \end{tikzpicture}
          \caption{方形开花弹示意图}
          \label{fig:square-danmaku}
        \end{figure}

        设图\ref{fig:square-danmaku}中
        方形$B_0B_8CD$的中心为原点, 边长为$2v_0$
        (使得方形开花弹以$v_0$的整体速度扩散),
        整体倾角为$\alpha$, $B_0,\dots,B_8$是一组间距相等的点.

        \begin{enumerate}[label=(\alph*)]
          \item 若倾角$\alpha=0\degree$,
                求点$B_i$ ($i=0,1,\dots,8$)
                的直角坐标$(v_{xi},v_{yi})$.
          \item 若倾角$\alpha\neq0\degree$,
                求点$B_i$ ($i=0,1,\dots,8$)
                的极坐标$\braket{v_i,\theta_i}$.
        \end{enumerate}

  \item \comment{激光与圆相交***}
        有一名弹设正在尝试新的弹幕样式,
        他设置了一些直线激光, 在自机处设置了一个圆形的``保护罩''.
        他希望当激光打到保护罩上时, 在激光与保护罩的交点生成子弹,
        同时应当把激光的终点设置到交点上 (通过控制激光的长度可以实现).
        具体的弹幕实现逻辑可能很复杂, 但简而言之他需要计算激光与圆的相交情况.

        图\ref{subfig:laser-circle-1}
        展示了他想实现的大致效果.
        在LuaSTG中, 我们容易获取激光的起点坐标, 朝向, 长度,
        以及圆形保护罩的圆心坐标, 半径,
        由这些信息足够确定激光与圆的相交情况.

        为了方便做进一步的计算,
        他以激光的起点为原点/极点, 激光朝向为$x$轴正方向/0\degree 方向,
        建立直角坐标系和极坐标系, 如图\ref{subfig:laser-circle-2}所示,
        $x$轴正半轴对应激光, 圆$O'$对应保护罩.

        \begin{figure}[htbp]
          \centering

          \begin{subfigure}{0.4\textwidth}
            \centering
            \begin{tikzpicture}[scale=3.5/384]
              \draw (-192,-224) rectangle (192,224);
              \fill[gray!50] (0,-150) circle (60);
              \draw[->] (-90,150) node[dot]{} -- +(-80:300-30);
            \end{tikzpicture}
            \caption{}
            \label{subfig:laser-circle-1}
          \end{subfigure}
          \begin{subfigure}{0.5\textwidth}
            \centering
            \begin{tikzpicture}[scale=1.5]
              \draw[-latex] (0,-0.5) -- (0,2) node[left] {$y$};
              \draw[-latex,name path=line]
              (-0.5,0) -- (4,0) node[right] {$x$};
              \draw[name path=circle,name=oo]
              (2,0.6) circle (0.9);

              \draw (0,0) node[dot,label={below left:$O$}]{}
              -- (2,0.6) node[dot,label={right:$O'$}]{};
              \draw (4mm,0) arc (0:15:4mm)
              node[above] {$\alpha$};

              \draw[name intersections={of=line and circle}]
              (intersection-1) node[dot,label={below left:$A$}]{}
              (intersection-2) node[dot,label={below right:$B$}]{};
            \end{tikzpicture}
            \caption{}
            \label{subfig:laser-circle-2}
          \end{subfigure}

          \caption{激光与圆相交}
        \end{figure}

        在新的坐标系下, 可以获取哪些信息呢? 这时激光的表示是非常简单的:
        起点为原点, 终点在$x$轴正半轴上, 长度等于终点的横坐标.
        保护罩圆心的极坐标比较容易获取:
        极径可以由$\Dist(激光,保护罩)$计算,
        极角可以由$(\Angle(激光,保护罩) - 激光朝向)$计算.

        设某一时刻, 保护罩的极坐标为$O'=\braket{L,\alpha}$,
        保护罩的半径为$R$. 请帮助他解决以下问题:

        \begin{enumerate}[label=(\alph*)]
          \item 若$L>R$, 那么$L,R,\alpha$满足什么条件时,
                圆$O'$与$x$轴\emph{正半轴}有交点?
          \item 若$L>R$, 且圆$O'$与$x$轴正半轴有$A,B$两个交点,
                且点$A$离原点更近,
                即图\ref{subfig:laser-circle-2}所示情况.
                求线段$OA$的长度.
          \item 若$L<R$, 那么圆$O'$与$x$轴正半轴一定有一个交点,
                设该交点为$A$. 求线段$OA$的长度.
        \end{enumerate}
\end{enumerate}
